\documentclass{article}
\usepackage[T1]{fontenc}
\usepackage{lmodern}
\usepackage{graphicx}
\usepackage{amsmath,amssymb}
\usepackage[a4paper,margin=2cm]{geometry}
\usepackage[absolute,overlay]{textpos}
\setlength{\TPHorizModule}{1mm}
\setlength{\TPVertModule}{1mm}
\usepackage[indonesian]{babel}
\usepackage{hyperref}
\setlength{\emergencystretch}{2em}
% unit: Halaman 34

\begin{document}

% === Halaman 34 (python/native) ===

(b) 2P = R




  m  ( 3(2)2 + 1)/ 8) (mod 11)  13/8 (mod 11)





               13  8–1 (mod 11)






  13  7 (mod 11) 





               78 mod 11  3 (mod 11)

Koordinat R:

  xr  m2 – 2xp  (mod p)  32 – 2  2  (mod 11)  5 (mod 11)

     yr  m(xp – xr) – yp (mod p)  3(2 – 5) – 4 (mod 11) 





  -13  (mod 11)  9 (mod 11)

 Jadi, R(5, 9)  






34

𝑚 3𝑥𝑝

2 + 𝑎

2𝑦𝑝

 (mod 𝑝)

\end{document}
